\documentclass[a4paper,12pt]{article}
\usepackage[utf8]{vietnam}
\usepackage{geometry}
\usepackage{enumitem}
\usepackage{titlesec}
\usepackage{array}
\usepackage{longtable}
\usepackage{booktabs}
\usepackage{tikz}
\usepackage{pgfplots}
\usepackage{graphicx}
\usetikzlibrary{shapes,arrows,positioning,fit,calc,backgrounds}
\pgfplotsset{compat=1.18}

\geometry{
    left=2.5cm,
    right=2.5cm,
    top=2.5cm,
    bottom=2.5cm
}

\titleformat{\section}
{\normalfont\Large\bfseries}{\thesection}{1em}{}
\titleformat{\subsection}
{\normalfont\large\bfseries}{\thesubsection}{1em}{}

\begin{document}

\begin{center}
\Large\textbf{ĐẶC TẢ USE CASE}\\
\large\textbf{Hệ thống Quản lý Kho}
\end{center}

\vspace{0.5cm}

% Bảng tổng quan Use Case
\section{Bảng tổng quan Use Case}
\begin{table}[h]
\centering
\renewcommand{\arraystretch}{1.5}
\begin{tabular}{|p{2cm}|p{3cm}|p{8cm}|}
\hline
\textbf{STT} & \textbf{Use Case} & \textbf{Mô tả chức năng} \\
\hline
3 & Quản lý nhà cung cấp & Lưu trữ và quản lý thông tin các nhà cung cấp, thêm, sửa, tìm kiếm nhà cung cấp \\
\hline
4 & Quản lý các kho & Quản lý thông tin các kho/khu vực lưu trữ, thêm, sửa thông tin kho, quản lý trạng thái kho \\
\hline
\end{tabular}
\caption{Tổng quan Use Case nghiệp vụ 3 và 4}
\end{table}

% Sơ đồ Use Case
\section{Sơ đồ Use Case}
\begin{figure}[h]
\centering
\begin{tikzpicture}[node distance=2cm, auto]
    % Define styles
    \tikzstyle{actor} = [rectangle, draw, fill=blue!20, text width=2.5cm, text centered, rounded corners, minimum height=1cm]
    \tikzstyle{usecase} = [ellipse, draw, fill=green!20, text width=3cm, text centered, minimum height=1cm]
    \tikzstyle{system} = [rectangle, draw, dashed, inner sep=10pt]
    \tikzstyle{line} = [draw, -latex']
    
    % Actors
    \node [actor] (admin) {Quản trị viên};
    \node [actor, below of=admin, node distance=3cm] (manager) {Quản lý kho};
    \node [actor, below of=manager, node distance=3cm] (staff) {Nhân viên kho};
    
    % System boundary
    \node [system, right of=admin, node distance=5cm, minimum width=8cm, minimum height=10cm] (system) {};
    \node [above of=system, node distance=5.5cm] at (system.north) {\textbf{Hệ thống Quản lý Kho}};
    
    % Use Cases for Supplier Management
    \node [usecase, right of=admin, node distance=6cm] (uc3_1) {Thêm nhà cung cấp};
    \node [usecase, below of=uc3_1, node distance=1.8cm] (uc3_2) {Sửa nhà cung cấp};
    \node [usecase, below of=uc3_2, node distance=1.8cm] (uc3_3) {Tìm kiếm nhà cung cấp};
    \node [usecase, below of=uc3_3, node distance=1.8cm] (uc3_4) {Xóa nhà cung cấp};
    
    % Use Cases for Warehouse Management
    \node [usecase, below of=uc3_4, node distance=1.8cm] (uc4_1) {Thêm kho};
    \node [usecase, below of=uc4_1, node distance=1.8cm] (uc4_2) {Sửa thông tin kho};
    \node [usecase, below of=uc4_2, node distance=1.8cm] (uc4_3) {Quản lý trạng thái kho};
    \node [usecase, below of=uc4_3, node distance=1.8cm] (uc4_4) {Xóa kho};
    \node [usecase, below of=uc4_4, node distance=1.8cm] (uc4_5) {Xem chi tiết kho};
    
    % Connections - Admin
    \path [line] (admin) -- (uc3_1);
    \path [line] (admin) -- (uc3_2);
    \path [line] (admin) -- (uc3_3);
    \path [line] (admin) -- (uc3_4);
    \path [line] (admin) -- (uc4_1);
    \path [line] (admin) -- (uc4_2);
    \path [line] (admin) -- (uc4_3);
    \path [line] (admin) -- (uc4_4);
    \path [line] (admin) -- (uc4_5);
    
    % Connections - Manager
    \path [line] (manager) -- (uc3_3);
    \path [line] (manager) -- (uc4_2);
    \path [line] (manager) -- (uc4_3);
    \path [line] (manager) -- (uc4_5);
    
    % Connections - Staff
    \path [line] (staff) -- (uc3_3);
    \path [line] (staff) -- (uc4_5);
    
    % Include relationships
    \path [line, dashed] (uc3_2) -- (uc3_3);
    \path [line, dashed] (uc4_2) -- (uc4_5);
\end{tikzpicture}
\caption{Sơ đồ Use Case cho nghiệp vụ Quản lý nhà cung cấp và Quản lý kho}
\end{figure}

\newpage

% Use Case 3: Quản lý nhà cung cấp
\section{Use Case 3: Quản lý nhà cung cấp}

\subsection{Mô tả}
Use case này cho phép người dùng lưu trữ và quản lý thông tin các nhà cung cấp, bao gồm thêm mới, chỉnh sửa thông tin và tìm kiếm nhà cung cấp trong hệ thống.

\subsection{Diễn viên (Actors)}
\begin{itemize}
    \item \textbf{Quản trị viên}: Người có quyền thêm, sửa, xóa và tìm kiếm thông tin nhà cung cấp
    \item \textbf{Nhân viên kho}: Người có quyền xem và tìm kiếm thông tin nhà cung cấp
\end{itemize}

\subsection{Điều kiện tiên quyết (Preconditions)}
\begin{itemize}
    \item Người dùng đã đăng nhập vào hệ thống
    \item Người dùng có quyền truy cập chức năng quản lý nhà cung cấp
    \item Hệ thống đang hoạt động bình thường
\end{itemize}

\subsection{Điều kiện hậu quả (Postconditions)}
\begin{itemize}
    \item Thông tin nhà cung cấp được lưu trữ thành công trong cơ sở dữ liệu
    \item Thông tin nhà cung cấp được cập nhật khi có thay đổi
    \item Danh sách nhà cung cấp được hiển thị chính xác
\end{itemize}

\subsection{Luồng chính (Main Flow)}

% Bảng chi tiết UC3.1
\begin{table}[h]
\centering
\renewcommand{\arraystretch}{1.3}
\begin{tabular}{|p{1.5cm}|p{10cm}|}
\hline
\textbf{Bước} & \textbf{Mô tả} \\
\hline
1 & Người dùng chọn chức năng "Quản lý nhà cung cấp" từ menu chính \\
\hline
2 & Hệ thống hiển thị danh sách các nhà cung cấp hiện có \\
\hline
3 & Người dùng nhấn nút "Thêm nhà cung cấp mới" \\
\hline
4 & Hệ thống hiển thị form nhập thông tin nhà cung cấp \\
\hline
5 & Người dùng nhập thông tin vào các trường \\
\hline
6 & Người dùng nhấn nút "Lưu" \\
\hline
7 & Hệ thống kiểm tra tính hợp lệ của dữ liệu \\
\hline
8 & Nếu dữ liệu hợp lệ, hệ thống lưu thông tin vào cơ sở dữ liệu \\
\hline
9 & Hệ thống hiển thị thông báo "Thêm nhà cung cấp thành công" \\
\hline
10 & Hệ thống cập nhật danh sách nhà cung cấp và quay lại màn hình danh sách \\
\hline
\end{tabular}
\caption{Luồng chính UC3.1: Thêm nhà cung cấp mới}
\end{table}

\textbf{UC3.1: Thêm nhà cung cấp mới}
\begin{enumerate}
    \item Người dùng chọn chức năng "Quản lý nhà cung cấp" từ menu chính
    \item Hệ thống hiển thị danh sách các nhà cung cấp hiện có
    \item Người dùng nhấn nút "Thêm nhà cung cấp mới"
    \item Hệ thống hiển thị form nhập thông tin nhà cung cấp với các trường:
    \begin{itemize}
        \item Mã nhà cung cấp (tự động hoặc nhập thủ công)
        \item Tên nhà cung cấp
        \item Địa chỉ
        \item Số điện thoại
        \item Email
        \item Người liên hệ
        \item Ghi chú
    \end{itemize}
    \item Người dùng nhập thông tin vào các trường
    \item Người dùng nhấn nút "Lưu"
    \item Hệ thống kiểm tra tính hợp lệ của dữ liệu:
    \begin{itemize}
        \item Kiểm tra các trường bắt buộc đã được nhập
        \item Kiểm tra định dạng email đúng
        \item Kiểm tra số điện thoại đúng định dạng
        \item Kiểm tra mã nhà cung cấp không trùng lặp
    \end{itemize}
    \item Nếu dữ liệu hợp lệ, hệ thống lưu thông tin vào cơ sở dữ liệu
    \item Hệ thống hiển thị thông báo "Thêm nhà cung cấp thành công"
    \item Hệ thống cập nhật danh sách nhà cung cấp và quay lại màn hình danh sách
\end{enumerate}

% Bảng chi tiết UC3.2
\begin{table}[h]
\centering
\renewcommand{\arraystretch}{1.3}
\begin{tabular}{|p{1.5cm}|p{10cm}|}
\hline
\textbf{Bước} & \textbf{Mô tả} \\
\hline
1 & Người dùng chọn chức năng "Quản lý nhà cung cấp" từ menu chính \\
\hline
2 & Hệ thống hiển thị danh sách các nhà cung cấp hiện có \\
\hline
3 & Người dùng chọn nhà cung cấp cần sửa từ danh sách \\
\hline
4 & Hệ thống hiển thị form chỉnh sửa thông tin nhà cung cấp với dữ liệu hiện tại \\
\hline
5 & Người dùng chỉnh sửa các thông tin cần thay đổi \\
\hline
6 & Người dùng nhấn nút "Cập nhật" \\
\hline
7 & Hệ thống kiểm tra tính hợp lệ của dữ liệu \\
\hline
8 & Nếu dữ liệu hợp lệ, hệ thống cập nhật thông tin trong cơ sở dữ liệu \\
\hline
9 & Hệ thống hiển thị thông báo "Cập nhật thông tin thành công" \\
\hline
10 & Hệ thống cập nhật danh sách nhà cung cấp và quay lại màn hình danh sách \\
\hline
\end{tabular}
\caption{Luồng chính UC3.2: Sửa thông tin nhà cung cấp}
\end{table}

\textbf{UC3.2: Sửa thông tin nhà cung cấp}
\begin{enumerate}
    \item Người dùng chọn chức năng "Quản lý nhà cung cấp" từ menu chính
    \item Hệ thống hiển thị danh sách các nhà cung cấp hiện có
    \item Người dùng chọn nhà cung cấp cần sửa từ danh sách
    \item Hệ thống hiển thị form chỉnh sửa thông tin nhà cung cấp với dữ liệu hiện tại
    \item Người dùng chỉnh sửa các thông tin cần thay đổi
    \item Người dùng nhấn nút "Cập nhật"
    \item Hệ thống kiểm tra tính hợp lệ của dữ liệu
    \item Nếu dữ liệu hợp lệ, hệ thống cập nhật thông tin trong cơ sở dữ liệu
    \item Hệ thống hiển thị thông báo "Cập nhật thông tin thành công"
    \item Hệ thống cập nhật danh sách nhà cung cấp và quay lại màn hình danh sách
\end{enumerate}

% Bảng chi tiết UC3.3
\begin{table}[h]
\centering
\renewcommand{\arraystretch}{1.3}
\begin{tabular}{|p{1.5cm}|p{10cm}|}
\hline
\textbf{Bước} & \textbf{Mô tả} \\
\hline
1 & Người dùng chọn chức năng "Quản lý nhà cung cấp" từ menu chính \\
\hline
2 & Hệ thống hiển thị danh sách các nhà cung cấp hiện có \\
\hline
3 & Người dùng nhập từ khóa tìm kiếm vào ô tìm kiếm \\
\hline
4 & Người dùng có thể chọn tiêu chí tìm kiếm (mã, tên, địa chỉ, số điện thoại) \\
\hline
5 & Người dùng nhấn nút "Tìm kiếm" \\
\hline
6 & Hệ thống tìm kiếm trong cơ sở dữ liệu theo tiêu chí đã chọn \\
\hline
7 & Hệ thống hiển thị danh sách kết quả tìm kiếm \\
\hline
8 & Nếu không tìm thấy kết quả, hệ thống hiển thị thông báo "Không tìm thấy nhà cung cấp phù hợp" \\
\hline
\end{tabular}
\caption{Luồng chính UC3.3: Tìm kiếm nhà cung cấp}
\end{table}

\textbf{UC3.3: Tìm kiếm nhà cung cấp}
\begin{enumerate}
    \item Người dùng chọn chức năng "Quản lý nhà cung cấp" từ menu chính
    \item Hệ thống hiển thị danh sách các nhà cung cấp hiện có
    \item Người dùng nhập từ khóa tìm kiếm vào ô tìm kiếm
    \item Người dùng có thể chọn tiêu chí tìm kiếm:
    \begin{itemize}
        \item Theo mã nhà cung cấp
        \item Theo tên nhà cung cấp
        \item Theo địa chỉ
        \item Theo số điện thoại
    \end{itemize}
    \item Người dùng nhấn nút "Tìm kiếm"
    \item Hệ thống tìm kiếm trong cơ sở dữ liệu theo tiêu chí đã chọn
    \item Hệ thống hiển thị danh sách kết quả tìm kiếm
    \item Nếu không tìm thấy kết quả, hệ thống hiển thị thông báo "Không tìm thấy nhà cung cấp phù hợp"
\end{enumerate}

% Bảng chi tiết UC3.4
\begin{table}[h]
\centering
\renewcommand{\arraystretch}{1.3}
\begin{tabular}{|p{1.5cm}|p{10cm}|}
\hline
\textbf{Bước} & \textbf{Mô tả} \\
\hline
1 & Người dùng chọn chức năng "Quản lý nhà cung cấp" từ menu chính \\
\hline
2 & Hệ thống hiển thị danh sách các nhà cung cấp hiện có \\
\hline
3 & Người dùng chọn nhà cung cấp cần xóa từ danh sách \\
\hline
4 & Hệ thống hiển thị thông báo xác nhận xóa \\
\hline
5 & Nếu người dùng chọn "Có", hệ thống kiểm tra xem nhà cung cấp có đang được sử dụng không \\
\hline
6 & Nếu không, hệ thống xóa nhà cung cấp khỏi cơ sở dữ liệu \\
\hline
7 & Hệ thống hiển thị thông báo "Xóa nhà cung cấp thành công" \\
\hline
8 & Nếu người dùng chọn "Không", hệ thống quay lại màn hình danh sách \\
\hline
\end{tabular}
\caption{Luồng chính UC3.4: Xóa nhà cung cấp}
\end{table}

\textbf{UC3.4: Xóa nhà cung cấp}
\begin{enumerate}
    \item Người dùng chọn chức năng "Quản lý nhà cung cấp" từ menu chính
    \item Hệ thống hiển thị danh sách các nhà cung cấp hiện có
    \item Người dùng chọn nhà cung cấp cần xóa từ danh sách
    \item Hệ thống hiển thị thông báo xác nhận "Bạn có chắc chắn muốn xóa nhà cung cấp này?"
    \item Nếu người dùng chọn "Có":
    \begin{itemize}
        \item Hệ thống kiểm tra xem nhà cung cấp có đang được sử dụng trong các phiếu nhập hàng không
        \item Nếu không, hệ thống xóa nhà cung cấp khỏi cơ sở dữ liệu
        \item Hệ thống hiển thị thông báo "Xóa nhà cung cấp thành công"
    \end{itemize}
    \item Nếu người dùng chọn "Không", hệ thống quay lại màn hình danh sách
\end{enumerate}

\subsection{Luồng thay thế (Alternative Flows)}

\textbf{UC3.1-A1: Dữ liệu không hợp lệ}
\begin{itemize}
    \item Tại bước 7 của UC3.1, nếu dữ liệu không hợp lệ:
    \item Hệ thống hiển thị thông báo lỗi cụ thể (ví dụ: "Email không đúng định dạng", "Số điện thoại không hợp lệ")
    \item Hệ thống giữ nguyên form nhập liệu để người dùng sửa lại
    \item Quay lại bước 5
\end{itemize}

\textbf{UC3.1-A2: Mã nhà cung cấp trùng lặp}
\begin{itemize}
    \item Tại bước 7 của UC3.1, nếu mã nhà cung cấp đã tồn tại:
    \item Hệ thống hiển thị thông báo "Mã nhà cung cấp đã tồn tại"
    \item Hệ thống yêu cầu người dùng nhập mã khác
    \item Quay lại bước 5
\end{itemize}

\textbf{UC3.4-A1: Nhà cung cấp đang được sử dụng}
\begin{itemize}
    \item Tại bước 5 của UC3.4, nếu nhà cung cấp đang được sử dụng trong phiếu nhập hàng:
    \item Hệ thống hiển thị thông báo "Không thể xóa nhà cung cấp này vì đang được sử dụng trong các phiếu nhập hàng"
    \item Hệ thống quay lại màn hình danh sách
\end{itemize}

\newpage

% Use Case 4: Quản lý các kho
\section{Use Case 4: Quản lý các kho}

\subsection{Mô tả}
Use case này cho phép người dùng quản lý thông tin các kho hoặc khu vực lưu trữ trong hệ thống, bao gồm thêm mới, chỉnh sửa thông tin kho và quản lý trạng thái hoạt động của kho.

\subsection{Diễn viên (Actors)}
\begin{itemize}
    \item \textbf{Quản trị viên}: Người có quyền thêm, sửa, xóa và quản lý trạng thái kho
    \item \textbf{Quản lý kho}: Người có quyền xem, sửa thông tin cơ bản và quản lý trạng thái kho
    \item \textbf{Nhân viên kho}: Người có quyền xem thông tin kho
\end{itemize}

\subsection{Điều kiện tiên quyết (Preconditions)}
\begin{itemize}
    \item Người dùng đã đăng nhập vào hệ thống
    \item Người dùng có quyền truy cập chức năng quản lý kho
    \item Hệ thống đang hoạt động bình thường
\end{itemize}

\subsection{Điều kiện hậu quả (Postconditions)}
\begin{itemize}
    \item Thông tin kho được lưu trữ thành công trong cơ sở dữ liệu
    \item Thông tin kho được cập nhật khi có thay đổi
    \item Trạng thái kho được cập nhật chính xác
    \item Danh sách kho được hiển thị chính xác
\end{itemize}

\subsection{Luồng chính (Main Flow)}

% Bảng chi tiết UC4.1
\begin{table}[h]
\centering
\renewcommand{\arraystretch}{1.3}
\begin{tabular}{|p{1.5cm}|p{10cm}|}
\hline
\textbf{Bước} & \textbf{Mô tả} \\
\hline
1 & Người dùng chọn chức năng "Quản lý kho" từ menu chính \\
\hline
2 & Hệ thống hiển thị danh sách các kho hiện có \\
\hline
3 & Người dùng nhấn nút "Thêm kho mới" \\
\hline
4 & Hệ thống hiển thị form nhập thông tin kho \\
\hline
5 & Người dùng nhập thông tin vào các trường \\
\hline
6 & Người dùng nhấn nút "Lưu" \\
\hline
7 & Hệ thống kiểm tra tính hợp lệ của dữ liệu \\
\hline
8 & Nếu dữ liệu hợp lệ, hệ thống lưu thông tin vào cơ sở dữ liệu \\
\hline
9 & Hệ thống thiết lập trạng thái mặc định cho kho là "Hoạt động" \\
\hline
10 & Hệ thống hiển thị thông báo "Thêm kho thành công" \\
\hline
11 & Hệ thống cập nhật danh sách kho và quay lại màn hình danh sách \\
\hline
\end{tabular}
\caption{Luồng chính UC4.1: Thêm kho mới}
\end{table}

\textbf{UC4.1: Thêm kho mới}
\begin{enumerate}
    \item Người dùng chọn chức năng "Quản lý kho" từ menu chính
    \item Hệ thống hiển thị danh sách các kho hiện có
    \item Người dùng nhấn nút "Thêm kho mới"
    \item Hệ thống hiển thị form nhập thông tin kho với các trường:
    \begin{itemize}
        \item Mã kho (tự động hoặc nhập thủ công)
        \item Tên kho
        \item Địa chỉ kho
        \item Diện tích kho
        \item Sức chứa tối đa
        \item Loại kho (kho nguyên liệu, kho thành phẩm, kho bán thành phẩm)
        \item Người phụ trách
        \item Số điện thoại liên hệ
        \item Ghi chú
    \end{itemize}
    \item Người dùng nhập thông tin vào các trường
    \item Người dùng nhấn nút "Lưu"
    \item Hệ thống kiểm tra tính hợp lệ của dữ liệu:
    \begin{itemize}
        \item Kiểm tra các trường bắt buộc đã được nhập
        \item Kiểm tra mã kho không trùng lặp
        \item Kiểm tra diện tích và sức chứa là số dương
    \end{itemize}
    \item Nếu dữ liệu hợp lệ, hệ thống lưu thông tin vào cơ sở dữ liệu
    \item Hệ thống thiết lập trạng thái mặc định cho kho là "Hoạt động"
    \item Hệ thống hiển thị thông báo "Thêm kho thành công"
    \item Hệ thống cập nhật danh sách kho và quay lại màn hình danh sách
\end{enumerate}

% Bảng chi tiết UC4.2
\begin{table}[h]
\centering
\renewcommand{\arraystretch}{1.3}
\begin{tabular}{|p{1.5cm}|p{10cm}|}
\hline
\textbf{Bước} & \textbf{Mô tả} \\
\hline
1 & Người dùng chọn chức năng "Quản lý kho" từ menu chính \\
\hline
2 & Hệ thống hiển thị danh sách các kho hiện có \\
\hline
3 & Người dùng chọn kho cần sửa từ danh sách \\
\hline
4 & Hệ thống hiển thị form chỉnh sửa thông tin kho với dữ liệu hiện tại \\
\hline
5 & Người dùng chỉnh sửa các thông tin cần thay đổi \\
\hline
6 & Người dùng nhấn nút "Cập nhật" \\
\hline
7 & Hệ thống kiểm tra tính hợp lệ của dữ liệu \\
\hline
8 & Nếu dữ liệu hợp lệ, hệ thống cập nhật thông tin trong cơ sở dữ liệu \\
\hline
9 & Hệ thống hiển thị thông báo "Cập nhật thông tin thành công" \\
\hline
10 & Hệ thống cập nhật danh sách kho và quay lại màn hình danh sách \\
\hline
\end{tabular}
\caption{Luồng chính UC4.2: Sửa thông tin kho}
\end{table}

\textbf{UC4.2: Sửa thông tin kho}
\begin{enumerate}
    \item Người dùng chọn chức năng "Quản lý kho" từ menu chính
    \item Hệ thống hiển thị danh sách các kho hiện có
    \item Người dùng chọn kho cần sửa từ danh sách
    \item Hệ thống hiển thị form chỉnh sửa thông tin kho với dữ liệu hiện tại
    \item Người dùng chỉnh sửa các thông tin cần thay đổi
    \item Người dùng nhấn nút "Cập nhật"
    \item Hệ thống kiểm tra tính hợp lệ của dữ liệu
    \item Nếu dữ liệu hợp lệ, hệ thống cập nhật thông tin trong cơ sở dữ liệu
    \item Hệ thống hiển thị thông báo "Cập nhật thông tin thành công"
    \item Hệ thống cập nhật danh sách kho và quay lại màn hình danh sách
\end{enumerate}

% Bảng chi tiết UC4.3
\begin{table}[h]
\centering
\renewcommand{\arraystretch}{1.3}
\begin{tabular}{|p{1.5cm}|p{10cm}|}
\hline
\textbf{Bước} & \textbf{Mô tả} \\
\hline
1 & Người dùng chọn chức năng "Quản lý kho" từ menu chính \\
\hline
2 & Hệ thống hiển thị danh sách các kho hiện có với trạng thái \\
\hline
3 & Người dùng chọn kho cần thay đổi trạng thái \\
\hline
4 & Người dùng chọn trạng thái mới từ danh sách \\
\hline
5 & Người dùng nhấn nút "Cập nhật trạng thái" \\
\hline
6 & Nếu trạng thái mới là "Đóng cửa" hoặc "Tạm dừng", hệ thống kiểm tra hàng hóa \\
\hline
7 & Hệ thống cập nhật trạng thái kho trong cơ sở dữ liệu \\
\hline
8 & Hệ thống hiển thị thông báo "Cập nhật trạng thái thành công" \\
\hline
9 & Hệ thống cập nhật danh sách kho \\
\hline
\end{tabular}
\caption{Luồng chính UC4.3: Quản lý trạng thái kho}
\end{table}

\textbf{UC4.3: Quản lý trạng thái kho}
\begin{enumerate}
    \item Người dùng chọn chức năng "Quản lý kho" từ menu chính
    \item Hệ thống hiển thị danh sách các kho hiện có với trạng thái
    \item Người dùng chọn kho cần thay đổi trạng thái
    \item Người dùng chọn trạng thái mới từ danh sách:
    \begin{itemize}
        \item Hoạt động
        \item Tạm dừng
        \item Đóng cửa
        \item Bảo trì
    \end{itemize}
    \item Người dùng nhấn nút "Cập nhật trạng thái"
    \item Nếu trạng thái mới là "Đóng cửa" hoặc "Tạm dừng":
    \begin{itemize}
        \item Hệ thống kiểm tra xem kho có hàng hóa đang lưu trữ không
        \item Nếu có, hệ thống yêu cầu người dùng xác nhận chuyển hàng sang kho khác
    \end{itemize}
    \item Hệ thống cập nhật trạng thái kho trong cơ sở dữ liệu
    \item Hệ thống hiển thị thông báo "Cập nhật trạng thái thành công"
    \item Hệ thống cập nhật danh sách kho
\end{enumerate}

% Bảng chi tiết UC4.4
\begin{table}[h]
\centering
\renewcommand{\arraystretch}{1.3}
\begin{tabular}{|p{1.5cm}|p{10cm}|}
\hline
\textbf{Bước} & \textbf{Mô tả} \\
\hline
1 & Người dùng chọn chức năng "Quản lý kho" từ menu chính \\
\hline
2 & Hệ thống hiển thị danh sách các kho hiện có \\
\hline
3 & Người dùng chọn kho cần xóa từ danh sách \\
\hline
4 & Hệ thống hiển thị thông báo xác nhận xóa \\
\hline
5 & Nếu người dùng chọn "Có", hệ thống kiểm tra hàng hóa và phiếu nhập/xuất \\
\hline
6 & Nếu không, hệ thống xóa kho khỏi cơ sở dữ liệu \\
\hline
7 & Hệ thống hiển thị thông báo "Xóa kho thành công" \\
\hline
8 & Nếu người dùng chọn "Không", hệ thống quay lại màn hình danh sách \\
\hline
\end{tabular}
\caption{Luồng chính UC4.4: Xóa kho}
\end{table}

\textbf{UC4.4: Xóa kho}
\begin{enumerate}
    \item Người dùng chọn chức năng "Quản lý kho" từ menu chính
    \item Hệ thống hiển thị danh sách các kho hiện có
    \item Người dùng chọn kho cần xóa từ danh sách
    \item Hệ thống hiển thị thông báo xác nhận "Bạn có chắc chắn muốn xóa kho này?"
    \item Nếu người dùng chọn "Có":
    \begin{itemize}
        \item Hệ thống kiểm tra xem kho có hàng hóa đang lưu trữ không
        \item Hệ thống kiểm tra xem kho có đang được sử dụng trong các phiếu nhập/xuất không
        \item Nếu không, hệ thống xóa kho khỏi cơ sở dữ liệu
        \item Hệ thống hiển thị thông báo "Xóa kho thành công"
    \end{itemize}
    \item Nếu người dùng chọn "Không", hệ thống quay lại màn hình danh sách
\end{enumerate}

% Bảng chi tiết UC4.5
\begin{table}[h]
\centering
\renewcommand{\arraystretch}{1.3}
\begin{tabular}{|p{1.5cm}|p{10cm}|}
\hline
\textbf{Bước} & \textbf{Mô tả} \\
\hline
1 & Người dùng chọn chức năng "Quản lý kho" từ menu chính \\
\hline
2 & Hệ thống hiển thị danh sách các kho hiện có \\
\hline
3 & Người dùng chọn kho cần xem chi tiết \\
\hline
4 & Hệ thống hiển thị thông tin chi tiết của kho \\
\hline
5 & Người dùng có thể chọn xem danh sách hàng hóa trong kho \\
\hline
\end{tabular}
\caption{Luồng chính UC4.5: Xem chi tiết kho}
\end{table}

\textbf{UC4.5: Xem chi tiết kho}
\begin{enumerate}
    \item Người dùng chọn chức năng "Quản lý kho" từ menu chính
    \item Hệ thống hiển thị danh sách các kho hiện có
    \item Người dùng chọn kho cần xem chi tiết
    \item Hệ thống hiển thị thông tin chi tiết của kho bao gồm:
    \begin{itemize}
        \item Thông tin cơ bản (mã, tên, địa chỉ, diện tích, sức chứa)
        \item Trạng thái hiện tại
        \item Số lượng hàng hóa đang lưu trữ
        \item Sức chứa còn lại
        \item Danh sách các khu vực trong kho
        \item Lịch sử hoạt động gần đây
    \end{itemize}
    \item Người dùng có thể chọn xem danh sách hàng hóa trong kho
\end{enumerate}

\subsection{Luồng thay thế (Alternative Flows)}

\textbf{UC4.1-A1: Dữ liệu không hợp lệ}
\begin{itemize}
    \item Tại bước 7 của UC4.1, nếu dữ liệu không hợp lệ:
    \item Hệ thống hiển thị thông báo lỗi cụ thể
    \item Hệ thống giữ nguyên form nhập liệu để người dùng sửa lại
    \item Quay lại bước 5
\end{itemize}

\textbf{UC4.1-A2: Mã kho trùng lặp}
\begin{itemize}
    \item Tại bước 7 của UC4.1, nếu mã kho đã tồn tại:
    \item Hệ thống hiển thị thông báo "Mã kho đã tồn tại"
    \item Hệ thống yêu cầu người dùng nhập mã khác
    \item Quay lại bước 5
\end{itemize}

\textbf{UC4.3-A1: Kho có hàng hóa khi đóng cửa}
\begin{itemize}
    \item Tại bước 6 của UC4.3, nếu kho có hàng hóa đang lưu trữ:
    \item Hệ thống hiển thị thông báo "Kho đang có hàng hóa, vui lòng chuyển hàng sang kho khác trước khi đóng cửa"
    \item Hệ thống hiển thị danh sách các kho hoạt động khác
    \item Người dùng phải chuyển hàng sang kho khác trước khi tiếp tục
\end{itemize}

\textbf{UC4.4-A1: Kho có hàng hóa hoặc đang được sử dụng}
\begin{itemize}
    \item Tại bước 5 của UC4.4, nếu kho có hàng hóa hoặc đang được sử dụng:
    \item Hệ thống hiển thị thông báo "Không thể xóa kho này vì đang có hàng hóa hoặc đang được sử dụng trong các phiếu nhập/xuất"
    \item Hệ thống quay lại màn hình danh sách
\end{itemize}

\newpage

% Sequence Diagram cho Ghi nhận hàng lỗi
\section{Sơ đồ trình tự (Sequence Diagram)}

\subsection{Sequence Diagram: Ghi nhận hàng lỗi}
\begin{figure}[h]
\centering
\begin{tikzpicture}[node distance=2.5cm, auto, >=stealth]
    % Define styles
    \tikzstyle{actor} = [rectangle, draw, fill=blue!20, text width=2cm, text centered, minimum height=0.8cm]
    \tikzstyle{lifeline} = [draw, dashed]
    \tikzstyle{message} = [draw, ->, thick]
    \tikzstyle{returnmsg} = [draw, ->, dashed, thick]
    \tikzstyle{altbox} = [draw, rectangle, fill=yellow!10, text width=12cm, minimum height=1cm]
    
    % Actors
    \node [actor] (nhanvien) {Nhân viên kho};
    \node [actor, right of=nhanvien, node distance=3.5cm] (giao_dien) {Giao diện quản lý hàng lỗi};
    \node [actor, right of=giao_dien, node distance=3.5cm] (controller) {HangLoiController};
    \node [actor, right of=controller, node distance=3.5cm] (database) {Database};
    
    % Lifelines
    \draw [lifeline] (nhanvien.south) -- ++(0,-10);
    \draw [lifeline] (giao_dien.south) -- ++(0,-10);
    \draw [lifeline] (controller.south) -- ++(0,-10);
    \draw [lifeline] (database.south) -- ++(0,-10);
    
    % Messages
    \node [text width=3cm, align=center] at (2.5, -0.5) {Chọn "Ghi nhận hàng lỗi"};
    \draw [message] (nhanvien.east) -- (giao_dien.west);
    
    \node [text width=3cm, align=center] at (6, -1) {ghiNhanHangLoi()};
    \draw [message] (giao_dien.east) -- (controller.west);
    
    \node [text width=3cm, align=center] at (9.5, -1.5) {Lấy danh sách hàng hóa};
    \draw [message] (controller.east) -- (database.west);
    
    \node [text width=3cm, align=center] at (9.5, -2) {Danh sách hàng hóa};
    \draw [returnmsg] (database.west) -- (controller.east);
    
    \node [text width=3cm, align=center] at (6, -2.5) {Danh sách hàng hóa};
    \draw [returnmsg] (controller.west) -- (giao_dien.east);
    
    \node [text width=3cm, align=center] at (2.5, -3) {Hiển thị danh sách};
    \draw [message] (giao_dien.west) -- (nhanvien.east);
    
    \node [text width=4cm, align=center] at (2.5, -3.8) {Chọn hàng hóa, nhập SL lỗi, chọn nguyên nhân};
    \draw [message] (nhanvien.east) -- (giao_dien.west);
    
    \node [text width=4cm, align=center] at (6, -4.3) {ghiNhanHangLoi(Hàng hóa, SL lỗi, nguyên nhân)};
    \draw [message] (giao_dien.east) -- (controller.west);
    
    % Alt box 1
    \draw [altbox] (3.5, -4.8) rectangle (11.5, -5.8);
    \node [text width=8cm, align=center] at (7.5, -5.3) {\textbf{alt} [Chưa chọn nguyên nhân]};
    \node [text width=3cm, align=center] at (6, -5.5) {"Vui lòng chọn nguyên nhân"};
    \draw [message] (controller.west) -- (giao_dien.east);
    
    % Alt box 2
    \draw [altbox] (3.5, -5.9) rectangle (11.5, -9.5);
    \node [text width=8cm, align=center] at (7.5, -6.3) {\textbf{alt} [Đã chọn nguyên nhân]};
    
    \node [text width=3cm, align=center] at (9.5, -6.8) {Kiểm tra số lượng tồn};
    \draw [message] (controller.east) -- (database.west);
    
    \node [text width=3cm, align=center] at (9.5, -7.3) {Số lượng tồn};
    \draw [returnmsg] (database.west) -- (controller.east);
    
    % Nested alt 2.1
    \draw [altbox, fill=red!10] (5, -7.6) rectangle (10, -8.3);
    \node [text width=4cm, align=center] at (7.5, -7.9) {[SL lỗi > SL tồn]};
    \node [text width=3cm, align=center] at (6, -8) {"Số lượng vượt quá tồn kho"};
    \draw [message] (controller.west) -- (giao_dien.east);
    
    % Nested alt 2.2
    \draw [altbox, fill=green!10] (5, -8.4) rectangle (10, -9.3);
    \node [text width=4cm, align=center] at (7.5, -8.7) {[Số lượng hợp lệ]};
    
    \node [text width=3cm, align=center] at (9.5, -8.9) {Giảm tồn kho, chuyển sang khu vực lỗi, lưu phiếu};
    \draw [message] (controller.east) -- (database.west);
    
    \node [text width=3cm, align=center] at (9.5, -9.1) {Lưu thành công};
    \draw [returnmsg] (database.west) -- (controller.east);
    
    \node [text width=3cm, align=center] at (6, -9.2) {"Thông báo thành công"};
    \draw [message] (controller.west) -- (giao_dien.east);
    
\end{tikzpicture}
\caption{Sequence Diagram: Ghi nhận hàng lỗi}
\end{figure}

\subsection{Sequence Diagram: Xử lý hàng lỗi}
\begin{figure}[h]
\centering
\begin{tikzpicture}[node distance=2.5cm, auto, >=stealth]
    % Define styles
    \tikzstyle{actor} = [rectangle, draw, fill=blue!20, text width=2cm, text centered, minimum height=0.8cm]
    \tikzstyle{lifeline} = [draw, dashed]
    \tikzstyle{message} = [draw, ->, thick]
    \tikzstyle{returnmsg} = [draw, ->, dashed, thick]
    \tikzstyle{altbox} = [draw, rectangle, fill=yellow!10, text width=12cm, minimum height=1cm]
    
    % Actors
    \node [actor] (quanly) {Quản lý kho};
    \node [actor, right of=quanly, node distance=3.5cm] (giao_dien) {Giao diện quản lý hàng lỗi};
    \node [actor, right of=giao_dien, node distance=3.5cm] (controller) {HangLoiController};
    \node [actor, right of=controller, node distance=3.5cm] (database) {Database};
    
    % Lifelines
    \draw [lifeline] (quanly.south) -- ++(0,-12);
    \draw [lifeline] (giao_dien.south) -- ++(0,-12);
    \draw [lifeline] (controller.south) -- ++(0,-12);
    \draw [lifeline] (database.south) -- ++(0,-12);
    
    % Messages
    \node [text width=3cm, align=center] at (2.5, -0.5) {Mở phiếu hàng lỗi};
    \draw [message] (quanly.east) -- (giao_dien.west);
    
    \node [text width=3cm, align=center] at (6, -1) {Yêu cầu mở phiếu (Mã phiếu)};
    \draw [message] (giao_dien.east) -- (controller.west);
    
    \node [text width=3cm, align=center] at (9.5, -1.5) {Lấy thông tin phiếu hàng lỗi};
    \draw [message] (controller.east) -- (database.west);
    
    \node [text width=3cm, align=center] at (9.5, -2) {Thông tin phiếu};
    \draw [returnmsg] (database.west) -- (controller.east);
    
    \node [text width=3cm, align=center] at (6, -2.5) {Thông tin phiếu};
    \draw [returnmsg] (controller.west) -- (giao_dien.east);
    
    % Alt box 1
    \draw [altbox] (3.5, -2.8) rectangle (11.5, -5.5);
    \node [text width=8cm, align=center] at (7.5, -3.2) {\textbf{alt}};
    
    % Alt 1.1
    \draw [altbox, fill=red!10] (4, -3.4) rectangle (11, -4.2);
    \node [text width=6cm, align=center] at (7.5, -3.7) {[Phiếu đã xử lý]};
    \node [text width=4cm, align=center] at (6, -3.9) {"Phiếu đã được xử lý, không thể xử lý lại"};
    \draw [message] (giao_dien.west) -- (quanly.east);
    
    % Alt 1.2
    \draw [altbox, fill=green!10] (4, -4.3) rectangle (11, -5.3);
    \node [text width=6cm, align=center] at (7.5, -4.6) {[Phiếu chưa xử lý]};
    
    \node [text width=3cm, align=center] at (2.5, -4.8) {Hiển thị thông tin phiếu};
    \draw [message] (giao_dien.west) -- (quanly.east);
    
    \node [text width=4cm, align=center] at (6, -5) {Hiển thị phương án xử lý (Trả nhà cung cấp, Tiêu hủy, Sửa chữa)};
    \draw [message] (giao_dien.west) -- (quanly.east);
    
    \node [text width=3cm, align=center] at (2.5, -5.8) {Chọn phương án xử lý};
    \draw [message] (quanly.east) -- (giao_dien.west);
    
    % Alt box 2
    \draw [altbox] (3.5, -6) rectangle (11.5, -8.5);
    \node [text width=8cm, align=center] at (7.5, -6.4) {\textbf{alt}};
    
    % Alt 2.1
    \draw [altbox, fill=red!10] (4, -6.6) rectangle (11, -7.2);
    \node [text width=6cm, align=center] at (7.5, -6.8) {[Chưa chọn phương án xử lý]};
    \node [text width=4cm, align=center] at (6, -7) {"Vui lòng chọn phương án xử lý"};
    \draw [message] (giao_dien.west) -- (quanly.east);
    
    % Alt 2.2
    \draw [altbox, fill=green!10] (4, -7.3) rectangle (11, -8.3);
    \node [text width=6cm, align=center] at (7.5, -7.6) {[Đã chọn phương án xử lý]};
    
    \node [text width=3cm, align=center] at (2.5, -7.8) {Xác nhận xử lý};
    \draw [message] (quanly.east) -- (giao_dien.west);
    
    \node [text width=3cm, align=center] at (6, -8) {Xử lý hàng lỗi (phương án)};
    \draw [message] (giao_dien.east) -- (controller.west);
    
    \node [text width=3cm, align=center] at (9.5, -8.2) {Cập nhật tồn kho hàng lỗi};
    \draw [message] (controller.east) -- (database.west);
    
    \node [text width=3cm, align=center] at (9.5, -8.4) {Lưu lịch sử xử lý};
    \draw [message] (controller.east) -- (database.west);
    
    \node [text width=3cm, align=center] at (9.5, -8.6) {Cập nhật trạng thái phiếu = Đã xử lý};
    \draw [message] (controller.east) -- (database.west);
    
    \node [text width=3cm, align=center] at (9.5, -8.8) {Cập nhật thành công};
    \draw [returnmsg] (database.west) -- (controller.east);
    
    \node [text width=3cm, align=center] at (6, -9) {Thông báo xử lý thành công};
    \draw [message] (controller.west) -- (giao_dien.east);
    
    \node [text width=3cm, align=center] at (2.5, -9.2) {Thông báo xử lý thành công};
    \draw [message] (giao_dien.west) -- (quanly.east);
    
\end{tikzpicture}
\caption{Sequence Diagram: Xử lý hàng lỗi}
\end{figure}

\newpage

% Sequence Diagram cho Quản lý nhà cung cấp
\subsection{Sequence Diagram: Thêm nhà cung cấp}
\begin{figure}[h]
\centering
\begin{tikzpicture}[node distance=2.5cm, auto, >=stealth]
    % Define styles
    \tikzstyle{actor} = [rectangle, draw, fill=blue!20, text width=2cm, text centered, minimum height=0.8cm]
    \tikzstyle{lifeline} = [draw, dashed]
    \tikzstyle{message} = [draw, ->, thick]
    \tikzstyle{returnmsg} = [draw, ->, dashed, thick]
    \tikzstyle{altbox} = [draw, rectangle, fill=yellow!10, text width=12cm, minimum height=1cm]
    
    % Actors
    \node [actor] (nguoidung) {Người dùng};
    \node [actor, right of=nguoidung, node distance=3.5cm] (giao_dien) {Giao diện quản lý NCC};
    \node [actor, right of=giao_dien, node distance=3.5cm] (controller) {NhaCungCapController};
    \node [actor, right of=controller, node distance=3.5cm] (database) {Database};
    
    % Lifelines
    \draw [lifeline] (nguoidung.south) -- ++(0,-8);
    \draw [lifeline] (giao_dien.south) -- ++(0,-8);
    \draw [lifeline] (controller.south) -- ++(0,-8);
    \draw [lifeline] (database.south) -- ++(0,-8);
    
    % Messages
    \node [text width=3cm, align=center] at (2.5, -0.5) {Chọn "Thêm NCC"};
    \draw [message] (nguoidung.east) -- (giao_dien.west);
    
    \node [text width=3cm, align=center] at (6, -1) {Hiển thị form thêm NCC};
    \draw [message] (giao_dien.west) -- (nguoidung.east);
    
    \node [text width=4cm, align=center] at (2.5, -1.8) {Nhập thông tin NCC};
    \draw [message] (nguoidung.east) -- (giao_dien.west);
    
    \node [text width=3cm, align=center] at (2.5, -2.5) {Nhấn "Lưu"};
    \draw [message] (nguoidung.east) -- (giao_dien.west);
    
    \node [text width=4cm, align=center] at (6, -3) {themNhaCungCap(thông tin)};
    \draw [message] (giao_dien.east) -- (controller.west);
    
    % Alt box
    \draw [altbox] (3.5, -3.3) rectangle (11.5, -7.5);
    \node [text width=8cm, align=center] at (7.5, -3.7) {\textbf{alt}};
    
    % Alt 1: Dữ liệu không hợp lệ
    \draw [altbox, fill=red!10] (4, -3.9) rectangle (11, -4.7);
    \node [text width=6cm, align=center] at (7.5, -4.2) {[Dữ liệu không hợp lệ]};
    \node [text width=4cm, align=center] at (6, -4.4) {"Thông báo lỗi cụ thể"};
    \draw [message] (controller.west) -- (giao_dien.east);
    
    % Alt 2: Mã trùng lặp
    \draw [altbox, fill=red!10] (4, -4.8) rectangle (11, -5.6);
    \node [text width=6cm, align=center] at (7.5, -5.1) {[Mã NCC trùng lặp]};
    \node [text width=4cm, align=center] at (6, -5.3) {"Mã NCC đã tồn tại"};
    \draw [message] (controller.west) -- (giao_dien.east);
    
    % Alt 3: Dữ liệu hợp lệ
    \draw [altbox, fill=green!10] (4, -5.7) rectangle (11, -7.3);
    \node [text width=6cm, align=center] at (7.5, -6) {[Dữ liệu hợp lệ]};
    
    \node [text width=3cm, align=center] at (9.5, -6.3) {Lưu thông tin NCC};
    \draw [message] (controller.east) -- (database.west);
    
    \node [text width=3cm, align=center] at (9.5, -6.6) {Lưu thành công};
    \draw [returnmsg] (database.west) -- (controller.east);
    
    \node [text width=3cm, align=center] at (6, -6.9) {"Thêm NCC thành công"};
    \draw [message] (controller.west) -- (giao_dien.east);
    
    \node [text width=3cm, align=center] at (2.5, -7.1) {"Thêm NCC thành công"};
    \draw [message] (giao_dien.west) -- (nguoidung.east);
    
\end{tikzpicture}
\caption{Sequence Diagram: Thêm nhà cung cấp}
\end{figure}

\subsection{Sequence Diagram: Tìm kiếm nhà cung cấp}
\begin{figure}[h]
\centering
\begin{tikzpicture}[node distance=2.5cm, auto, >=stealth]
    % Define styles
    \tikzstyle{actor} = [rectangle, draw, fill=blue!20, text width=2cm, text centered, minimum height=0.8cm]
    \tikzstyle{lifeline} = [draw, dashed]
    \tikzstyle{message} = [draw, ->, thick]
    \tikzstyle{returnmsg} = [draw, ->, dashed, thick]
    
    % Actors
    \node [actor] (nguoidung) {Người dùng};
    \node [actor, right of=nguoidung, node distance=3.5cm] (giao_dien) {Giao diện quản lý NCC};
    \node [actor, right of=giao_dien, node distance=3.5cm] (controller) {NhaCungCapController};
    \node [actor, right of=controller, node distance=3.5cm] (database) {Database};
    
    % Lifelines
    \draw [lifeline] (nguoidung.south) -- ++(0,-5);
    \draw [lifeline] (giao_dien.south) -- ++(0,-5);
    \draw [lifeline] (controller.south) -- ++(0,-5);
    \draw [lifeline] (database.south) -- ++(0,-5);
    
    % Messages
    \node [text width=3cm, align=center] at (2.5, -0.5) {Chọn "Tìm kiếm NCC"};
    \draw [message] (nguoidung.east) -- (giao_dien.west);
    
    \node [text width=3cm, align=center] at (6, -1) {Hiển thị ô tìm kiếm};
    \draw [message] (giao_dien.west) -- (nguoidung.east);
    
    \node [text width=4cm, align=center] at (2.5, -1.8) {Nhập từ khóa, chọn tiêu chí};
    \draw [message] (nguoidung.east) -- (giao_dien.west);
    
    \node [text width=3cm, align=center] at (2.5, -2.5) {Nhấn "Tìm kiếm"};
    \draw [message] (nguoidung.east) -- (giao_dien.west);
    
    \node [text width=4cm, align=center] at (6, -3) {timKiemNhaCungCap(từ khóa, tiêu chí)};
    \draw [message] (giao_dien.east) -- (controller.west);
    
    \node [text width=3cm, align=center] at (9.5, -3.5) {Tìm kiếm theo tiêu chí};
    \draw [message] (controller.east) -- (database.west);
    
    \node [text width=3cm, align=center] at (9.5, -4) {Kết quả tìm kiếm};
    \draw [returnmsg] (database.west) -- (controller.east);
    
    \node [text width=3cm, align=center] at (6, -4.5) {Kết quả tìm kiếm};
    \draw [returnmsg] (controller.west) -- (giao_dien.east);
    
    \node [text width=3cm, align=center] at (2.5, -4.8) {Hiển thị kết quả};
    \draw [message] (giao_dien.west) -- (nguoidung.east);
    
\end{tikzpicture}
\caption{Sequence Diagram: Tìm kiếm nhà cung cấp}
\end{figure}

\newpage

% Sequence Diagram cho Quản lý các kho
\subsection{Sequence Diagram: Thêm kho}
\begin{figure}[h]
\centering
\begin{tikzpicture}[node distance=2.5cm, auto, >=stealth]
    % Define styles
    \tikzstyle{actor} = [rectangle, draw, fill=blue!20, text width=2cm, text centered, minimum height=0.8cm]
    \tikzstyle{lifeline} = [draw, dashed]
    \tikzstyle{message} = [draw, ->, thick]
    \tikzstyle{returnmsg} = [draw, ->, dashed, thick]
    \tikzstyle{altbox} = [draw, rectangle, fill=yellow!10, text width=12cm, minimum height=1cm]
    
    % Actors
    \node [actor] (nguoidung) {Người dùng};
    \node [actor, right of=nguoidung, node distance=3.5cm] (giao_dien) {Giao diện quản lý kho};
    \node [actor, right of=giao_dien, node distance=3.5cm] (controller) {KhoController};
    \node [actor, right of=controller, node distance=3.5cm] (database) {Database};
    
    % Lifelines
    \draw [lifeline] (nguoidung.south) -- ++(0,-8);
    \draw [lifeline] (giao_dien.south) -- ++(0,-8);
    \draw [lifeline] (controller.south) -- ++(0,-8);
    \draw [lifeline] (database.south) -- ++(0,-8);
    
    % Messages
    \node [text width=3cm, align=center] at (2.5, -0.5) {Chọn "Thêm kho"};
    \draw [message] (nguoidung.east) -- (giao_dien.west);
    
    \node [text width=3cm, align=center] at (6, -1) {Hiển thị form thêm kho};
    \draw [message] (giao_dien.west) -- (nguoidung.east);
    
    \node [text width=4cm, align=center] at (2.5, -1.8) {Nhập thông tin kho};
    \draw [message] (nguoidung.east) -- (giao_dien.west);
    
    \node [text width=3cm, align=center] at (2.5, -2.5) {Nhấn "Lưu"};
    \draw [message] (nguoidung.east) -- (giao_dien.west);
    
    \node [text width=4cm, align=center] at (6, -3) {themKho(thông tin)};
    \draw [message] (giao_dien.east) -- (controller.west);
    
    % Alt box
    \draw [altbox] (3.5, -3.3) rectangle (11.5, -7.5);
    \node [text width=8cm, align=center] at (7.5, -3.7) {\textbf{alt}};
    
    % Alt 1: Dữ liệu không hợp lệ
    \draw [altbox, fill=red!10] (4, -3.9) rectangle (11, -4.7);
    \node [text width=6cm, align=center] at (7.5, -4.2) {[Dữ liệu không hợp lệ]};
    \node [text width=4cm, align=center] at (6, -4.4) {"Thông báo lỗi cụ thể"};
    \draw [message] (controller.west) -- (giao_dien.east);
    
    % Alt 2: Mã trùng lặp
    \draw [altbox, fill=red!10] (4, -4.8) rectangle (11, -5.6);
    \node [text width=6cm, align=center] at (7.5, -5.1) {[Mã kho trùng lặp]};
    \node [text width=4cm, align=center] at (6, -5.3) {"Mã kho đã tồn tại"};
    \draw [message] (controller.west) -- (giao_dien.east);
    
    % Alt 3: Dữ liệu hợp lệ
    \draw [altbox, fill=green!10] (4, -5.7) rectangle (11, -7.3);
    \node [text width=6cm, align=center] at (7.5, -6) {[Dữ liệu hợp lệ]};
    
    \node [text width=3cm, align=center] at (9.5, -6.3) {Lưu thông tin kho, trạng thái = Hoạt động};
    \draw [message] (controller.east) -- (database.west);
    
    \node [text width=3cm, align=center] at (9.5, -6.6) {Lưu thành công};
    \draw [returnmsg] (database.west) -- (controller.east);
    
    \node [text width=3cm, align=center] at (6, -6.9) {"Thêm kho thành công"};
    \draw [message] (controller.west) -- (giao_dien.east);
    
    \node [text width=3cm, align=center] at (2.5, -7.1) {"Thêm kho thành công"};
    \draw [message] (giao_dien.west) -- (nguoidung.east);
    
\end{tikzpicture}
\caption{Sequence Diagram: Thêm kho}
\end{figure}

\subsection{Sequence Diagram: Quản lý trạng thái kho}
\begin{figure}[h]
\centering
\begin{tikzpicture}[node distance=2.5cm, auto, >=stealth]
    % Define styles
    \tikzstyle{actor} = [rectangle, draw, fill=blue!20, text width=2cm, text centered, minimum height=0.8cm]
    \tikzstyle{lifeline} = [draw, dashed]
    \tikzstyle{message} = [draw, ->, thick]
    \tikzstyle{returnmsg} = [draw, ->, dashed, thick]
    \tikzstyle{altbox} = [draw, rectangle, fill=yellow!10, text width=12cm, minimum height=1cm]
    
    % Actors
    \node [actor] (nguoidung) {Người dùng};
    \node [actor, right of=nguoidung, node distance=3.5cm] (giao_dien) {Giao diện quản lý kho};
    \node [actor, right of=giao_dien, node distance=3.5cm] (controller) {KhoController};
    \node [actor, right of=controller, node distance=3.5cm] (database) {Database};
    
    % Lifelines
    \draw [lifeline] (nguoidung.south) -- ++(0,-8);
    \draw [lifeline] (giao_dien.south) -- ++(0,-8);
    \draw [lifeline] (controller.south) -- ++(0,-8);
    \draw [lifeline] (database.south) -- ++(0,-8);
    
    % Messages
    \node [text width=3cm, align=center] at (2.5, -0.5) {Chọn kho cần thay đổi};
    \draw [message] (nguoidung.east) -- (giao_dien.west);
    
    \node [text width=3cm, align=center] at (6, -1) {Hiển thị trạng thái hiện tại};
    \draw [message] (giao_dien.west) -- (nguoidung.east);
    
    \node [text width=3cm, align=center] at (2.5, -1.8) {Chọn trạng thái mới};
    \draw [message] (nguoidung.east) -- (giao_dien.west);
    
    \node [text width=3cm, align=center] at (2.5, -2.5) {Nhấn "Cập nhật trạng thái"};
    \draw [message] (nguoidung.east) -- (giao_dien.west);
    
    \node [text width=4cm, align=center] at (6, -3) {capNhatTrangThaiKho(mã kho, trạng thái)};
    \draw [message] (giao_dien.east) -- (controller.west);
    
    % Alt box
    \draw [altbox] (3.5, -3.3) rectangle (11.5, -7.5);
    \node [text width=8cm, align=center] at (7.5, -3.7) {\textbf{alt}};
    
    % Alt 1: Đóng cửa/Tạm dừng có hàng hóa
    \draw [altbox, fill=red!10] (4, -3.9) rectangle (11, -5.2);
    \node [text width=6cm, align=center] at (7.5, -4.2) {[Đóng cửa/Tạm dừng có hàng hóa]};
    
    \node [text width=3cm, align=center] at (9.5, -4.5) {Kiểm tra hàng hóa trong kho};
    \draw [message] (controller.east) -- (database.west);
    
    \node [text width=3cm, align=center] at (9.5, -4.8) {Có hàng hóa};
    \draw [returnmsg] (database.west) -- (controller.east);
    
    \node [text width=4cm, align=center] at (6, -5) {"Kho có hàng hóa, vui lòng chuyển hàng"};
    \draw [message] (controller.west) -- (giao_dien.east);
    
    % Alt 2: Cập nhật thành công
    \draw [altbox, fill=green!10] (4, -5.3) rectangle (11, -7.3);
    \node [text width=6cm, align=center] at (7.5, -5.6) {[Cập nhật thành công]};
    
    \node [text width=3cm, align=center] at (9.5, -5.9) {Cập nhật trạng thái kho};
    \draw [message] (controller.east) -- (database.west);
    
    \node [text width=3cm, align=center] at (9.5, -6.2) {Cập nhật thành công};
    \draw [returnmsg] (database.west) -- (controller.east);
    
    \node [text width=3cm, align=center] at (6, -6.5) {"Cập nhật trạng thái thành công"};
    \draw [message] (controller.west) -- (giao_dien.east);
    
    \node [text width=3cm, align=center] at (2.5, -6.8) {"Cập nhật trạng thái thành công"};
    \draw [message] (giao_dien.west) -- (nguoidung.east);
    
\end{tikzpicture}
\caption{Sequence Diagram: Quản lý trạng thái kho}
\end{figure}

\end{document}
